监督学习里有一列标签当裁判，模型好不好，拿它一比就知道。聚类没有这个东西。数据不会主动告诉你哪种分组是对的，于是每个聚类算法开工前都得先自带一个裁判——把``好的分组''写成一个可以计算的准则。K-means 写的是``每个点尽量贴近所属簇的欧氏中心''，谱聚类写的是``簇内由高权重路径连通、簇间连接稀薄''，层次聚类则把判断交给 linkage 规则，由它决定哪些局部合并值得保留。

不同的准则给出不同的分组，而且这些分组常常都说得通。

\begin{center}
\begin{tikzpicture}[>=stealth]
% ---- 左：分成两组 ----
\draw[thick] (0,0.1) rectangle (3.4,3.4);
\foreach \p in {(0.60,0.70),(1.15,0.62),(0.80,1.20),(1.22,1.08),(0.95,0.85),(0.62,1.02),
                (0.60,2.30),(1.15,2.22),(0.80,2.80),(1.22,2.68),(0.95,2.45),(0.62,2.62),
                (2.20,0.70),(2.75,0.62),(2.40,1.20),(2.82,1.08),(2.55,0.85),(2.22,1.02),
                (2.20,2.30),(2.75,2.22),(2.40,2.80),(2.82,2.68),(2.55,2.45),(2.22,2.62)}
  \fill \p circle (2.2pt);
\draw[very thick] (1.7,0.15) -- (1.7,3.35);
\node at (1.7,-0.45) {\footnotesize 两组：左列对右列};
% ---- 右：分成四组 ----
\begin{scope}[shift={(4.6,0)}]
\draw[thick] (0,0.1) rectangle (3.4,3.4);
\foreach \p in {(0.60,0.70),(1.15,0.62),(0.80,1.20),(1.22,1.08),(0.95,0.85),(0.62,1.02),
                (0.60,2.30),(1.15,2.22),(0.80,2.80),(1.22,2.68),(0.95,2.45),(0.62,2.62),
                (2.20,0.70),(2.75,0.62),(2.40,1.20),(2.82,1.08),(2.55,0.85),(2.22,1.02),
                (2.20,2.30),(2.75,2.22),(2.40,2.80),(2.82,2.68),(2.55,2.45),(2.22,2.62)}
  \fill \p circle (2.2pt);
\draw[very thick] (1.7,0.15) -- (1.7,3.35);
\draw[very thick] (0.05,1.7) -- (3.35,1.7);
\node at (1.7,-0.45) {\footnotesize 四组：每团各算一个};
\end{scope}
\end{tikzpicture}
\end{center}

\emph{同一批点，两种分法都讲得通，数据里没有裁判来判谁对。}

图里那两种切法不是一对一错，而是同一份数据在两个尺度上的结构。既然如此，``评价''就不能被安排在算法之后当作验收环节——它已经参与了问题的定义。你选定的目标函数即是你对簇的定义，改目标就是改问题，两次运行的结果不可比，因为它们回答的不是同一个问题。

本单元从 K-means 的平方距离目标和交替下降讲起，把三件容易被混为一谈的事情分开：目标单调下降是算法性质，停在局部最优是非凸目标的性质，而分出来的簇有没有现实含义则完全是第三件事。随后把样本变成加权图，用 Laplacian 二次型、Rayleigh 商和连续松弛解释谱聚类为何能识别弯曲的、非凸的簇，以及它把尺度选择挪到了哪里。全程反复出现的追问只有几个：距离怎么定义、簇数怎么定、结果换个种子还在不在。

PCA、流形学习与矩阵补全留在前一单元，因为它们先回答的是表示与重构的问题。本单元只在需要时调用这些表示，同时提醒一件容易被忽略的事：先降维再聚类，被聚类的对象已经换成了数据在新表示下的像，问题本身被改写过了。高维下这个诱惑格外大，因为距离集中会让原始空间里的簇结构变得极脆——目标值在不同分区之间的差距缩进噪声量级，换个随机种子就能翻盘。

\subsection{怎么读这一章}

第一遍先用 K-means 建立一个认识：簇是目标函数与几何共同的产物，而不是数据里等着被发现的实体。然后用谱聚类做对照实验，观察换一张相似图之后，同一批点如何长出另一套分区。特征值松弛的完整推导可以暂时略过，尺度参数、簇数、初始化和稳定性检查不能略过——绝大多数聚类事故出在这四处，而不是出在算法实现上。

一个自查的标准是：如果你说不出换一种度量或换一套邻接规则为什么会换来另一种``真相''，就先别给簇起名字。名字一旦落进报表，那个簇就获得了它并不具备的实体感。

\subsection{本章篇目}

以下篇目按概念依赖排列；若从具体问题进入，也可以先打开最接近当前困惑的一篇，再沿篇内链接回补。

\begin{itemize}
\tightlist
\item \hyperref[sec:13-Machine-Learning/05-Clustering-and-Data-Geometry/01]{``K-means 在交替下降中制造簇''}；
\item \hyperref[sec:13-Machine-Learning/05-Clustering-and-Data-Geometry/02]{``谱聚类把连通性变成特征向量''}。
\end{itemize}

两篇读完，你会有两套互相对照的簇定义，以及一份在动手前就该填的表：用了哪把尺子、建了哪张图、定了几个簇、换个种子还剩多少。这份交代做完，聚类结果才谈得上可以被别人检验。
